\section{The software pipeline}
\label{sec:pipeline}

The package turns one TOML file into one results directory. Each stage is a module with a
single job, so each can be tested and replaced on its own. Table~\ref{tab:pipeline} lists them in
the order a run visits them.

\begin{table}[H]
  \centering
  \small
  \begin{tabular}{L{3.2cm}L{4.6cm}L{7.4cm}}
    \toprule
    stage & module & role \\
    \midrule
    configuration & \py{config.py} & pydantic models for lattice, ansatz, optimiser, backend; validated TOML in \py{configs/} \\
    lattice & \py{lattice.py} & immutable bond list and \py{site_order} (\cref{dfn:ring}), independent of Qiskit \\
    Hamiltonian & \py{hamiltonian.py} & \py{SparsePauliOp} of \cref{dfn:ham}; exact reference by sparse Lanczos \\
    ansatz & \py{ansatz/} & \py{RVBAnsatz}, \py{HVAnsatz}, library baselines; \py{unshare.py} for parameter shift (\cref{sec:m1}) \\
    compilation & \py{transpile.py} & preset pass manager (level 3), \py{initial_layout} on the heavy-hex loop, \py{observable.apply_layout} \\
    estimation & \py{estimators.py} & exact statevector, sampled Aer (ideal or noisy), IBM Runtime; one protocol \py{EnergyEstimator} \\
    optimisation & \py{optimizers/} & in-repo ADAM, SPSA, gradient descent; gradient strategies (forward, central, SPSA, parameter shift) \\
    driver & \py{vqe.py} & objective with trace, ideal-energy tracking, final re-measurement \\
    persistence & \py{results.py} & \py{results/<UTC>-<name>/}: config, result, traces, environment, log \\
    \bottomrule
  \end{tabular}
  \caption{Pipeline stages of \texttt{kagome-vqe} (\py{src/kagome_vqe/}).}
  \label{tab:pipeline}
\end{table}

\subsection{Compilation for the device}

The target is \py{FakeGuadalupeV2}, Qiskit's snapshot of the 16-qubit \texttt{ibmq\_guadalupe}
including its calibrated noise. Its coupling map is heavy-hex, with a 12-qubit loop
$(1,2,3,5,8,11,14,13,12,10,7,4)$ (\py{GUADALUPE_HEAVY_HEX_LOOP}). Passing this loop as the initial
layout puts every ring bond on a physical coupling, so the ring gates need no routing SWAPs. The
observable then follows the circuit through \py{apply_layout}; this is the explicit Qiskit~2
replacement for the accident described in \cref{sec:ansatz}.

\begin{table}[H]
  \centering
  % Generated by `uv run python -m docsgen`; do not edit.
\begin{tabular}{lcccccc}
\toprule
ansatz & parameters & gate occurrences & PS energies & CX (Guadalupe) & CX (unshared) & depth \\
\midrule
HVA & \VParamsHva & $2\times 18$ & \VPsEnergiesHva & \VCxHva & \VCxHvaUnshared & \VDepthHva \\
RVB & \VParamsRvb & $12$ & \VPsEnergiesRvb & \VCxRvb & \VCxRvbUnshared & \VDepthRvb \\
\bottomrule
\end{tabular}

  \caption{Circuit cost on \py{FakeGuadalupeV2} (optimisation level 3, heavy-hex loop layout).
    ``PS energies'' is the number of energies one parameter-shift gradient needs
    (\cref{sec:m1}). ``Unshared'' is the circuit with one parameter per gate occurrence,
    which parameter shift requires; it has the same two-qubit cost.}
  \label{tab:circuits}
\end{table}

Both ansatzes cost \VCxHva{} CX after transpilation, and the count does not depend on the
layout hint (\VCxRvbNoLayout{} for RVB without it), so routing is not the cause. The count
decomposes as $6+12\times5$: six singlet CX and five CX for each of the 12 Heisenberg bond
gates. A general two-qubit unitary needs at most three CX, so the bond gates are not being
resynthesised.

\begin{prop}[Binding parameters before transpiling]\label{prop:bound}
  The circuits are transpiled once with \emph{unbound} parameters. Level-3 two-qubit
  resynthesis (block consolidation followed by unitary synthesis) needs numeric matrices and
  skips parametric gates, so each $R_{ZZ}R_{YY}R_{XX}$ block keeps its naive decomposition
  minus cancellations. Binding the parameters first gives $\VBoundCxHva=6+12\times3$ CX (depth
  \VBoundDepthHva{} instead of \VDepthHva{} for HVA). At $\theta_k=\VBoundTheta$ for every
  parameter, under the full \py{FakeGuadalupeV2} noise model, the fraction of the exact energy
  that survives rises from \VParamKeepHva{} to \VBoundKeepHva{} (HVA) and from \VParamKeepRvb{}
  to \VBoundKeepRvb{} (RVB).
\end{prop}

\begin{proof}
  Measured, not derived: \py{docsgen.physics.bound_compilation} transpiles both versions with
  the same layout and seed and evaluates each with \VNoiseShots{} shots per basis. The CX
  counts are the transpiler's output. The comparison uses $\theta\neq0$ because at zero the
  bound circuit collapses to the 6-CX singlet preparation, which would be a trivial win
  (\cref{rem:easy}).
\end{proof}

\begin{rem}[Consequence for the roadmap]
  Re-transpiling for every parameter vector costs classical time for each energy (seconds, at
  level 3), which matters for parameter shift with \VPsEnergiesHva{} energies per step. Two cheaper
  routes give the same 3-CX blocks: transpile the bond gate once as the textbook 3-CX
  decomposition of $\exp(-i(aXX+bYY+cZZ))$ with parametric single-qubit rotations, or bind and
  transpile only the 12 distinct two-qubit blocks. Either one is the first, nearly free step of
  roadmap M4 (depth compression), and it should raise the noisy floor before any approximate
  method is tried.
\end{rem}

\begin{rem}[A stale figure in the prose docs]
  Until \texttt{687d0ae}, \py{docs/physics.md} (section~8) said the transpiled circuits have
  ``66 (HVA) to 93 (RVB)'' CX, a sentence from the initial commit \texttt{d4a05ac}. With the
  current code RVB compiles to \VCxRvb{} CX with or without the layout (\cref{tab:circuits}),
  equal to HVA, as \cref{lem:eswap} predicts; the legacy eSWAP gives 54. The corrected text
  gives 66 for both and attributes 93 to the 2022 notebook's layout; see \cref{sec:status}.
\end{rem}

\subsection{Estimating an energy}

The Hamiltonian splits into three groups of qubit-wise commuting terms (all $XX$, all $YY$, all
$ZZ$), so one energy costs \VMeasGroups{} measured circuits. With $s$ shots per basis, the
estimator returns the sampled mean of each Pauli term and their weighted sum. Throughout,
``shots per energy'' means shots per measurement basis, and ``shots spent'' counts all bases.

\begin{dfn}[Relative error]
  For an energy $E$ and the exact $E_0$, $\varepsilon(E)=|E-E_0|/|E_0|$ (\py{relative_error}).
  The challenge asked for $\varepsilon<1\,\%$.
\end{dfn}

The package distinguishes three energies of a run, and the distinction matters from M0 on:
\begin{itemize}
  \item the \emph{trace} energies, one per evaluation the optimiser asked for (noisy, finite shots);
  \item the \emph{final} energy: the final parameters re-measured once with a large, fixed budget
    (\py{backend.final_shots}, 16384 per basis);
  \item the \emph{ideal} energy: the exact statevector energy of the current parameters, recorded
    after every iteration but never shown to the optimiser (\py{track_ideal_energy}). It measures
    how good the \emph{parameters} are, separately from the noise floor.
\end{itemize}
